\section{Introduction}
\label{sec:intro}

Each section of this paper lives in its own file in the \texttt{sections} folder, and \texttt{main.tex} reads them
with \verb|\input|. The outline on the left shows every section of every file, numbered as in the PDF, with its word
count.

Double-click anywhere in the PDF to jump to the line that made it. Press Ctrl+J to see where the cursor is in the
PDF. Hover a citation such as \citet{knuth1984texbook} to see the entry, or a reference such as
Section~\ref{sec:method} to see what it points at. Ctrl+click follows it.

Type \verb|\cite{| and the editor offers the entries of \texttt{references.bib}. Type \verb|\ref{| and it offers
your labels. To add a reference, open the command palette (Ctrl+Shift+P) and choose \emph{Add a reference from a DOI
or arXiv ID}. Classic introductions to the language are \citet{lamport1994latex} and \citet{mittelbach2004companion}.
